shall choose different indeterminates for each factor in the product of @@M0@@ with itself, and indicate the indeterminate as a superscript when necessary: @@M1@@, @@M2@@, etc. For @@M3@@, we therefore have @@M4@@. To define the operations @@M5@@ on @@M6@@, we shall therefore specify a morphism of functors: @@M7@@.
For this, we must first define a filtration of the functor @@M8@@ that permits the gluing of 1.6.

We first observe that @@M9@@ is an ideal for the law @@M10@@; for this, functoriality reduces us to the case @@M11@@ and the product @@M12@@, and the assertion follows from 2.3.2. We then take as @@M13@@ the set of series of the form @@M14@@, @@M15@@, where @@M16@@ denotes the integer part of @@M17@@.

The element defining the morphism @@M18@@ is @@M19@@; the reader will verify as an exercise that the conditions of 1.6 are indeed satisfied (in fact, the filtration is chosen ad hoc).

Here again, the operations @@M20@@ are calculated by substitution from @@M21@@, and are therefore determined by universal polynomials. We denote by @@M22@@ the coefficient of @@M23@@ in @@M24@@; it is an isobaric polynomial of weight @@M25@@ in the variables @@M26@@ (a verification analogous to that of 1.15). It is calculated by symmetric functions from the defining element:
\sgaStableTarget{sga6.e05.eq.2-3-3}@@M27@@

As an example, let us give the value of @@M28@@:
